% ============================================================
% §3  Methods
% ============================================================
\section{Methods}
\label{sec:methods}

\subsection{Bayesian Optimization Protocol}
\label{sec:protocol}

I follow the standard sequential BO protocol.
An initial design of $n_{\mathrm{init}} = 2d$ points is drawn via Sobol
quasi-random sampling, where $d$ is the problem dimensionality; Sobol
provides more uniform space-filling coverage than independent uniform
random draws.
After the initial design I run $T = 20$ sequential optimization
iterations.
At each step the GP surrogate is refitted to all accumulated observations,
the acquisition function is maximized to select one candidate point, and
the true objective is evaluated at that point.
All tensor computations use \texttt{torch.float64} to preserve GP
numerical stability.

\paragraph{Performance metric.}
Performance is measured by \emph{simple regret}
\[
  r_t = f^* - \max_{i \le t}\, y_i,
\]
where $f^*$ is the known optimal value and $y_i$ are observed function
values.
Because all six problems are framed as \emph{maximization},
$r_t \ge 0$ throughout every run and decreases toward zero as
the optimizer improves.

\paragraph{Experimental protocol.}
Each strategy–problem pair is evaluated across 10 independent random
seeds.
Acquisition functions are maximized using BoTorch's
\texttt{optimize\_acqf} with multi-start optimization
(\texttt{num\_restarts}\,$= 10$, \texttt{raw\_samples}\,$= 64$).

% ----------------------------------------------------------
\subsection{Gaussian Process Surrogate}
\label{sec:gp}

I use BoTorch's \texttt{SingleTaskGP}~\citep{balandat2020botorch} as
the surrogate: a Gaussian process with a radial basis function (RBF)
kernel equipped with automatic relevance determination (ARD), one
independent lengthscale per input dimension.
Kernel hyperparameters—lengthscales, output scale, and noise
variance—are fitted by maximizing the exact marginal log-likelihood
(MLL) at every iteration.

\paragraph{Input normalization.}
A critical design choice is to normalize GP inputs to the unit
hypercube $[0,1]^d$ and standardize outputs to zero mean and unit
variance before fitting.
I apply BoTorch's \texttt{Normalize(d,\,bounds=problem.bounds)} and
\texttt{Standardize(m=1)} transforms, using the \emph{known, fixed}
problem bounds rather than data-adaptive bounds.
Fixing the normalization range to the full search space prevents it
from drifting as observations accumulate, giving the GP a stable
length-scale calibration throughout the run.

The necessity of this choice was revealed empirically by the two
engineering problems.
Borehole and Piston have inputs spanning approximately seven orders of
magnitude (e.g.\ Piston's initial gas volume
$V_0 \in [0.002,\,0.010]\,\mathrm{m}^3$ alongside ambient pressure
$P_0 \in [90\,000,\,110\,000]\,\mathrm{Pa}$, a factor of roughly
$10^7$).
Without normalization, the GP cannot learn sensible per-dimension
lengthscales across these scales, causing EI and UCB to degrade to
random-search performance.
Table~\ref{tab:normalization} shows the effect on Piston, where the
controlled comparison is cleanest: adding normalization reduces EI's
mean final regret from $0.41$ to $0.019$ (${\approx}20{\times}$) and
UCB's from $0.39$ to $0.0019$ (${\approx}200{\times}$), while the
Random baseline—which makes no GP calls—remains unchanged at
${\approx}0.40$.
The unchanged Random performance isolates the fault squarely to the
surrogate, not the problem formulation.

\begin{table}[h]
\centering
\caption{Effect of GP input normalization on mean final regret on
  Piston (7D), 10 seeds.  Random is unchanged because it does not use
  the GP, confirming the fault is in the surrogate.}
\label{tab:normalization}
\begin{tabular}{lccc}
\toprule
Strategy & Before norm. & After norm. & Improvement \\
\midrule
EI       & 0.41  & \textbf{0.019}  & ${\approx}20{\times}$  \\
UCB      & 0.39  & \textbf{0.0019} & ${\approx}200{\times}$ \\
Random   & 0.40  & 0.41            & —                      \\
\bottomrule
\end{tabular}
\end{table}

% ----------------------------------------------------------
\subsection{Acquisition Strategies}
\label{sec:strategies}

I benchmark four strategies that together span the full
exploration–exploitation spectrum (Table~\ref{tab:strategies}).

\paragraph{Expected Improvement (EI).}
\[
  \alpha_{\mathrm{EI}}(x)
    = \mathbb{E}\!\left[\max\!\bigl(f(x) - f^+,\, 0\bigr)\right],
\]
where $f^+ = \max_i y_i$ is the current best observed value.
Under a GP surrogate this has the closed form
$\alpha_{\mathrm{EI}}(x) =
  \sigma(x)\!\left[\gamma(x)\Phi(\gamma(x)) + \phi(\gamma(x))\right]$,
with $\gamma(x) = (\mu(x) - f^+)/\sigma(x)$, where $\Phi$ and $\phi$
are the standard Gaussian CDF and PDF.
EI rewards both a high posterior mean (exploitation) and a high
posterior variance (exploration); neither component alone is
sufficient to generate positive EI.
Implemented via \texttt{botorch.acquisition.ExpectedImprovement} with
\texttt{best\_f = train\_y.max()}.

\paragraph{Upper Confidence Bound (UCB).}
\[
  \alpha_{\mathrm{UCB}}(x) = \mu(x) + \sqrt{\beta}\,\sigma(x),
  \quad \beta = 2.
\]
The $\beta$ parameter controls the exploration–exploitation
trade-off.
Unlike EI, GP-UCB has finite-time regret guarantees under appropriate
$\beta$ schedules~\citep{frazier2018tutorial}.
Implemented via \texttt{botorch.acquisition.UpperConfidenceBound}.

\paragraph{Uncertainty (pure exploration).}
\[
  \alpha_{\mathrm{Unc}}(x) = \sigma(x).
\]
This strategy selects the point of maximum posterior standard
deviation, entirely ignoring the posterior mean.
It is the pure-exploration extreme of the spectrum: it uses the same
GP as EI and UCB but discards all exploitation signal.
Its inclusion allows me to isolate whether the GP's uncertainty
estimate \emph{alone}—without any mean information—is sufficient to
outperform blind search.
Implemented via \texttt{botorch.acquisition.PosteriorStandardDeviation}.

\paragraph{Random (baseline).}
Points are drawn from a Sobol sequence within the search bounds at
every iteration; the GP is never called.
This is the weakest baseline that any data-efficient strategy must
surpass.

\begin{table}[h]
\centering
\caption{The four acquisition strategies and their position on the
  exploration–exploitation spectrum.}
\label{tab:strategies}
\begin{tabular}{llcc}
\toprule
Strategy & Acquisition & Uses $\mu$? & Uses $\sigma$? \\
\midrule
EI          & ExpectedImprovement  & \checkmark             & \checkmark \\
UCB         & UpperConfidenceBound & \checkmark             & \checkmark \\
Uncertainty & PosteriorStdDev      & \phantom{\checkmark}   & \checkmark \\
Random      & Sobol sampling       & \phantom{\checkmark}   & \phantom{\checkmark} \\
\bottomrule
\end{tabular}
\end{table}

% ----------------------------------------------------------
\subsection{Problem Suite}
\label{sec:problems}

I benchmark on six functions spanning dimensions $d \in \{2,6,7,8,10\}$
across two categories (Table~\ref{tab:problems}).

\begin{table*}[t]
\centering
\caption{Benchmark problem suite.  All problems are framed as
  maximization; simple regret is $f^* - \text{best observed}$.
  Engineering optimal values are derived from exhaustive corner search
  ($2^d$ vertices) rather than random sampling, which systematically
  undershoots vertex maxima on monotone functions (see
  Section~\ref{sec:gp}).}
\label{tab:problems}
\begin{tabular}{llcll}
\toprule
Problem & Type & $d$ & Domain & $f^*$ \\
\midrule
Branin          & Synthetic    & 2  & standard          & $-0.3979$ \\
Six-Hump Camel  & Synthetic    & 2  & $[-3,3]\times[-2,2]$ & $1.0316$ \\
Hartmann-6      & Synthetic    & 6  & $[0,1]^6$         & $3.3224$  \\
Piston          & Engineering  & 7  & multi-scale       & $1.20$    \\
Borehole        & Engineering  & 8  & multi-scale       & $310.0$   \\
Ackley-10D      & Synthetic    & 10 & $[-5,5]^{10}$     & $0.0$     \\
\bottomrule
\end{tabular}
\end{table*}

\textbf{Synthetic problems} (Branin, Six-Hump Camel, Hartmann-6,
Ackley-10D) wrap BoTorch's built-in test functions with
\texttt{negate=True} to convert minimization to maximization;
optimal values are analytically known.
The Ackley bounds are narrowed from BoTorch's default
$[-32.768, 32.768]^{10}$ to the literature-standard $[-5, 5]^{10}$.

\textbf{Engineering problems} model physical systems with inputs at
physically meaningful, heterogeneous scales.
\emph{Borehole} (8D) models water flow through a borehole; inputs
include borehole radius, hydraulic transmissivities, water-table
elevations, and aquifer conductivity, sourced from SMT's
\texttt{WaterFlow}.
\emph{Piston} (7D) models piston cycle time as a function of mass $M$,
surface area $S$, initial gas volume $V_0$, spring stiffness $k$,
atmospheric pressure $P_0$, ambient temperature $T_a$, and gas
temperature $T_0$ \citep{surjanovic2013virtual}.
Piston is absent from the SMT version used, so it is implemented
manually from the standard closed-form formula.
For both engineering functions the global maximum lies at a box vertex;
optimal values are obtained by exhaustive corner evaluation
($2^8 = 256$ and $2^7 = 128$ vertices respectively) and a small
conservative buffer to guarantee non-negative simple regret.

% ----------------------------------------------------------
\subsection{Statistical Methodology}
\label{sec:stats}

With four strategies and six problems, I require a comparison
framework that is robust to non-Gaussian, heavy-tailed regret
distributions and avoids the inflated type-I error of repeated
pairwise $t$-tests.
Following \citet{demsar2006statistical}, I use the Friedman rank test
with Nemenyi post-hoc pairwise comparisons.

\paragraph{Friedman test.}
Each (problem, seed) pair forms one block; the four strategies are
ranked by final regret within each block, rank 1 indicating best.
This yields $N = 60$ blocks (6 problems $\times$ 10 seeds) and
$k = 4$ strategies.
The Friedman statistic tests the null hypothesis that all strategies
have equal expected rank.

\paragraph{Nemenyi post-hoc test.}
On rejection of the Friedman null, all $\binom{4}{2} = 6$ pairwise
comparisons are performed using Nemenyi's test at $\alpha = 0.05$
\citep{terpilowski2019scikit}.
Two strategies are deemed \emph{statistically indistinguishable} when
their average-rank difference is less than the critical difference
\[
  \mathrm{CD}
    = q_\alpha \sqrt{\frac{k(k+1)}{6N}}
    = q_{0.05}\sqrt{\frac{4 \times 5}{6 \times 60}},
\]
where $q_\alpha$ is the critical value from the Studentized range
distribution (divided by $\sqrt{2}$); $\mathrm{CD} \approx 0.61$ for
$k=4$, $N=60$, $\alpha=0.05$.
Results are visualized as a Critical Difference (CD) diagram
(Figure~\ref{fig:cd}), where horizontal bars connect strategies that
are not significantly different from one another.
The CD value is computed analytically inside the plotting function to
remain correct under any change to $N$ or $k$.
